\documentclass[12pt,a4paper]{article}
\usepackage[utf8]{inputenc}
\usepackage[margin=1in]{geometry}
\usepackage{graphicx}
\usepackage{xcolor}
\usepackage{tcolorbox}
\usepackage{tabularx}
\usepackage{booktabs}
\usepackage{enumitem}
\usepackage{hyperref}
\hypersetup{
    colorlinks=true,
    linkcolor=black,
    filecolor=black,
    urlcolor=black,
    citecolor=black
}

\usepackage{listings}

% Define colors
\definecolor{problemred}{RGB}{220,53,69}
\definecolor{solutiongreen}{RGB}{40,167,69}
\definecolor{hciblue}{RGB}{0,123,255}

\title{\textbf{Google Meet Redesign \& HCI Feature Mapping}\\
\large Final HCI Portfolio: Cognitive Comfort, Flexibility, and Error Prevention}
\author{CHALAMCHALA ANJANI NITHIN}
\date{\today}

\begin{document}

\begin{titlepage}
\begin{center}

{\LARGE \textbf{Indian Institute of Information Technology,\\
Design and Manufacturing, Kancheepuram}}\\[1.5cm]

{\large (HCI End-Semester Final Portfolio)}\\[3cm]

\rule{\linewidth}{0.5pt}\\[0.4cm]
{\Huge \textbf{The Google Meet Redesign}}\\[0.3cm]
{\Large \textbf{Cognitive Defense, Deep Heuristic Mapping \& Prototype Realization}}\\[0.3cm]
\rule{\linewidth}{0.5pt}\\[2cm]

\vfill

\textbf{CHALAMCHALA ANJANI NITHIN}\\[0.2cm]
\textbf{CS23B1102}\\[0.2cm]
\textit{Department of Computer Science and Engineering}\\[0.5cm]

{\large \today}

\end{center}
\end{titlepage}

\newpage
\tableofcontents
\newpage

\section{Introduction}

The objective of this final HCI portfolio is to present the complete architectural redesign of the Google Meet platform. Moving beyond traditional "prototype blocks", this report demonstrates fully realized HTML/CSS high-fidelity layouts, providing strict, in-depth mappings between the visual layout and fundamental cognitive design laws, Shneiderman's 8 Golden Rules, Nielsen's Usability Heuristics, and modern Interaction Design Paradigms.

This document specifically details how the application fulfills requirements ranging from the \textbf{Inverted Pyramid Pattern} to \textbf{Asimov's Laws of Error Prevention}.

\section{Enhanced Midsem features}
Since the midterm evaluation, the prototype has undergone extensive ``Phase II" enhancements to explicitly map against advanced psychological and classical design laws:

\begin{enumerate}
    \item \textbf{Inverted Pyramid Meeting Info Module}: A new layout for distributing critical session details hierarchically.
    \item \textbf{Robustness \& Connection Recovery}: Explicit user support systems during simulated faults.
    \item \textbf{Shape Psychology Alignment}: Re-structuring CSS tokens to enforce pill-shaped safe actions and sharp-cornered alerts.
    \item \textbf{Navigation Design Guidelines}: Introduction of explicit spatial breadcrumbs in deep settings menus.
\end{enumerate}

\section{Classical Design \& Cognitive Laws}

\subsection{The Vital Few (80-20 Rule) \& Miller's Law}
By stripping the primary control bar to exactly \textbf{7 primary buttons} (Mic, Camera, Screen Share, Hand Raise, Chat, Participants, Leave), we drastically lower the decision-making index, heavily leaning on the \textbf{Vital Few} theory where 80\% of meeting interaction occurs via these primary actions.

\subsection{Serial Position Effect (Primacy \& Recency)}
The positioning of the toolbar directly exploits human memory retention. 
\begin{itemize}
    \item \textbf{Primacy}: The Microphone and Video toggles are placed at the absolute start (left-side) of the toolbar, as they are the first actions processed visually.
    \item \textbf{Recency}: The highly-destructive "Leave" button is pinned to the extreme far-right, ensuring it is the final item recalled and avoiding adjacent misclicks.
\end{itemize}

\subsection{The Inverted Pyramid Information Structure}
Meeting Details originally dumped flat text onto the screen. Our redesign implements a strict \textbf{Inverted Pyramid} structure in the "Meeting Info Panel":
\begin{itemize}
    \item \textbf{Tier 1 (The Lead):} The absolutely crucial joining link and passcode, presented at the top with primary highlight colors.
    \item \textbf{Tier 2 (Body):} Helpful but secondary contexts (Host Name, Active Participant count).
    \item \textbf{Tier 3 (Tail):} General meeting etiquette and recording disclaimers.
\end{itemize}

\begin{figure}[h]
    \centering
    \begin{tcolorbox}[colback=gray!10,colframe=gray!50,arc=4mm,height=5cm,valign=center,halign=center]
        \large\textit{[ IMAGE PLACEHOLDER: Inverted Pyramid ]}\\[0.5cm]
        \small Open index.html $\rightarrow$ Click 'Info' in the control bar.
    \end{tcolorbox}
    \caption{The Inverted Pyramid applied to the Info Panel.}
\end{figure}

\section{Error Prevention, Learnability, \& Robustness}

\subsection{Asimov's Laws \& Robust System Status}
In UI design, Asimov's first law roughly translates to "A system must not allow the user to irreparably harm their session." To guarantee \textbf{Robustness}, we simulated a network failure state. Instead of dropping the call, the system renders a highly visible "Reconnecting..." safety overlay, temporarily disabling interactions to prevent chaotic data states.

\begin{figure}[h]
    \centering
    \begin{tcolorbox}[colback=gray!10,colframe=gray!50,arc=4mm,height=5cm,valign=center,halign=center]
        \large\textit{[ IMAGE PLACEHOLDER: Robustness Reconnect Mode ]}\\[0.5cm]
        \small Open index.html $\rightarrow$ Press Ctrl+Alt+N to simulate network failure.
    \end{tcolorbox}
    \caption{Robustness implementation handling system-level faults safely.}
\end{figure}

\subsection{Learnability (The Law of Learning)}
The \textbf{Law of Learning} states that proficiency is built via repetition and multi-modal reinforcement. Our Meeting UI tracks user clicks. If a user clicks the "Mute" button five times manually with their mouse, the system surfaces a contextual tooltip instructing them to use the \texttt{Spacebar} shortcut, gently pulling them up the power-law learning curve.

\section{Navigation \& Interaction Design Paradigms}

\subsection{Navigation Design Guidelines (Breadcrumbs)}
To satisfy modern navigation architecture guidelines, deeply nested Modals (e.g., Settings) now feature top-level \textbf{Breadcrumb trails} (Meeting $>$ Settings $>$ Audio). This provides immediate spatial context and ensures users never feel ``lost" inside the application state.

\subsection{Interaction Paradigms \& Mental Models}
We abandoned legacy text-CAPTCHAs in favor of a \textbf{Direct Manipulation Interaction Paradigm}. The authentication screen forces users to drag-and-drop a component into a slot. This spatial map perfectly mirrors the physical \textbf{Mental Model} of sliding a key into a lock, dramatically lowering cognitive load compared to deciphering distorted text.

\section{Visual Design Elements}

\subsection{Shape \& Color Psychology}
Color mapped the semantic states perfectly (Red = Danger, Blue = Safety). However, we mapped \textbf{Shape Psychology} natively into the CSS tokens:
\begin{itemize}
    \item \textbf{Fully Rounded / Pill Shapes}: Applied to primary, safe interaction nodes (the "Join" button). The lack of sharp edges psychologically implies approachability, community, and low risk.
    \item \textbf{Sharp Edges}: Applied exclusively to the destructive "Leave" button. The tighter radius functions as a visual subconscious stop-sign.
\end{itemize}

\section{Universal Accessibility (WCAG)}
The prototype integrates native physical HCI models for universal design:
\begin{itemize}
    \item \textbf{Motor-Assist Mode (Fitts's Law):} Tremors lead to erroneous disconnects on 48px targets. Motor-assist pads targets by 150\%.
    \item \textbf{Color-Blind High Contrast Matrix:} Extreme Yellow-on-Black tokens to survive Deuteranopia constraints.
\end{itemize}

\section{Conclusion}
This dynamic prototype fundamentally abandons legacy interactions. Through meticulous mapping of Classical Design Laws (Serial Position, Inverted Pyramid), Interaction Paradigms (Gestalt Rooms, Drag-CAPCHTA), and Psychological Triggers (Scarcity timers, Shape signaling, Nudges), the redesigned architecture proves that Web Usability Guidelines are not restrictive, but rather serve as the foundation for highly flexible, empathetic software.

\end{document}
